\section{Ethics Considerations}
\label{sec:ethics}


We took several measures to address ethical considerations in both the user survey and post-deletion behavioral analysis.
%
\textbf{(1) IRB review and informed consent.}
%
Our institution's Institutional Review Board (IRB) reviewed the survey and determined that it qualified as exempt human-subjects research (categories 2.i and 3A) on June 25, 2026.
%
We recruited U.S. adults through Prolific who had used at least one of the six LLM agents.
%
Participants provided informed consent before beginning the survey, were informed of its procedures and voluntary nature, and received \$4 upon completion, subject to the study's attention-check requirements.
%
\textbf{(2) Participant privacy and data protection.}
%
We did not collect participants' names, email addresses, or other direct identifiers, and used pseudonymous participant IDs provided by Prolific.
%
Survey data were stored on access-restricted, university-managed OneDrive, and the consent form permitted future research use and publication of de-identified data.
%
We also redacted potentially identifying information from participant-written deletion requests before using them in the post-deletion experiments.
%
\textbf{(3) Responsible testing of LLM agents.}
%
We conducted all experiments using researcher-controlled consumer accounts and synthetic information, avoiding the disclosure of real individuals' personal data.
%
Testing was limited to consumer-facing functions, including conversations, file uploads, memory settings, and deletion controls, without accessing other users' data or bypassing authentication or authorization mechanisms.
%
For two agents (XX and XX), automated sessions required a browser-stealth extension to bypass bot-detection checks.
%
This modified browser presentation but did not expand account permissions or grant access to protected resources.
%
Although we did not obtain explicit permission from the providers, we confined the experiments to our own accounts and synthetic data to minimize potential privacy and security risks.
%
\textbf{(4) Potential impacts on LLM agent providers.}
%
We recognize that reporting agent-specific findings on data deletion may raise reputational concerns for the evaluated providers.
%
However, we believe identifying such differences serves the public interest by improving transparency and accountability in data deletion practices.
%
To avoid misrepresenting the findings, we distinguish observable post-deletion recall from actual backend data retention or erasure and make no claims regarding legal noncompliance.
%
We believe these reporting safeguards, together with the limited scope of our experiments, help balance potential reputational risks against the public benefits of the research.